% 表 1：Individual Cognition & Choice
\begin{table*}[htbp]
    \centering
    \caption{Study summary and validation checks (\textbf{Category: Individual Cognition \& Choice})}
    \label{tab:study_summary_individual}
    \begin{tabular}{p{1cm}p{5.5cm}p{4cm}p{4cm}}
        \toprule
        \textbf{ID} & \textbf{Study Summary} & \textbf{Phenomenon-Level Tests} & \textbf{Data-Level Tests} \\
        \midrule
        
           \textbf{001} 
           & \textit{False Consensus Effect}~\cite{}: People overestimate how common their own choices are. \newline Replication of 3 sub-studies using diverse methods: hypothetical scenarios, personal trait questionnaires, and behavioral choices (sign experiment). N=320.
           & P1-P3: Sub-studies overall False Consensus Effect is statistically significant 
           & D1-D3:  Match False Consensus Effect magnitude (mean estimate difference) in sub-studies (e.g., sub-study in hypothetical scenarios: 17.2\%).  \\
           \cmidrule{1-4}
          
         \textbf{002} 
         & \textit{Anchoring Effect}~\cite{}: Initial values (anchors) influence subsequent numerical estimates. \newline General knowledge estimation task across 15 topics (e.g., Mississippi River length), comparing estimates under High vs. Low anchors. N=145. 
         & P1: There is a statistically significant anchoring effect \newline P2: High vs. low anchor produce different effects \newline P3: Anchoring inversely relates to confidence 
          & D1: Anchoring Index should match 0.47 \newline D2: Confidence should match 3.85 \\
         \cmidrule{1-4}
        
         \textbf{003} 
          & \textit{Framing Effect}~\cite{}: Identical choices elicit different preferences when framed as gains versus losses. \newline Implementation of 11 decision scenarios from the original paper, N=600 (between-subjects design). Scenarios are paired to test framing effects across gains/losses, certainty, and accounting.
          & P1: Shift from risk aversion to risk seeking. \newline P2: Preference for a sure win over a probabilistic win. \newline P3: Willingness to buy is higher after losing cash than after losing a ticket. \newline P4: Willingness to travel is higher when the base price is low than when it is high. %\newline Note: Prob 3/4 (dominance logic) and Prob 6 are implemented but not scored.
         & D1-D8: Match choice proportions across 8 problems (e.g., Asian Disease: 72\% risk averse in Gain vs. 78\% risk seeking in Loss). \\
          \cmidrule{1-4}
        
         \textbf{004} 
         & \textit{Representativeness Heuristic}~\cite{}: People judge probability by similarity to stereotypes rather than statistical principles. \newline Replication of 9 classic problems (e.g., Birth Sequence) demonstrating neglect of base rates and sample size. N=1500 (between-subjects design). 
          & P1-P8: Representativeness bias is present for each problem type (two of the 9 problems are combined into one between-subjects comparison test: P8)
          & D1: Birth sequence bias proportion should match 81.5\% \newline D2: Program choice bias proportion should match 75.3\% (only 2 problems have complete human baseline data) \\
        
        \bottomrule
    \end{tabular}
    
    \vspace{0.3cm}
    \small
    \textit{Note}: \textbf{Phenomenon-Level Tests (P)} verify whether the psychological effect is present. \textbf{Data-Level Tests (D)} measure how closely the quantitative results match human baseline data using equivalence testing with tolerance $\delta = 0.2$.
\end{table*}

% 表 2：Strategic Interaction
\begin{table*}[htbp]
    \centering
    \caption{Study summary and validation checks (\textbf{Category: Strategic Interaction})}
    \label{tab:study_summary_strategic}
    \begin{tabular}{p{1cm}p{5.5cm}p{4cm}p{3.5cm}}
        \toprule
        \textbf{ID} & \textbf{Study Summary} & \textbf{Phenomenon-Level Tests} & \textbf{Data-Level Tests} \\
        \midrule
        
        \textbf{005} & (TBD) & & \\ \cmidrule{1-4}
        \textbf{006} & (TBD) & & \\ \cmidrule{1-4}
        \textbf{007} & (TBD) & & \\ \cmidrule{1-4}
        \textbf{008} & (TBD) & & \\
        
        \bottomrule
    \end{tabular}
    
    \vspace{0.3cm}
    \small
%    \textit{Note}: \textbf{Phenomenon-Level Tests (P)} verify whether the psychological effect is present (critical for validation—must pass). \textbf{Data-Level Tests (D)} measure how closely the quantitative results match human baseline data using equivalence testing with tolerance $\delta = 0.2$ (non-critical—provides additional validation).
\end{table*}


% 表 3：Social Psychology
\begin{table*}[htbp]
    \centering
    \caption{Study summary and validation checks (\textbf{Category: Social Psychology})}
    \label{tab:study_summary_social}
    \begin{tabular}{p{1cm}p{5.5cm}p{4cm}p{3.5cm}}
        \toprule
        \textbf{Study} & \textbf{Study Summary} & \textbf{Phenomenon-Level Tests} & \textbf{Data-Level Tests} \\
        \midrule
        
        \textbf{009} & (TBD) & & \\ \cmidrule{1-4}
        \textbf{010} & (TBD) & & \\ \cmidrule{1-4}
        \textbf{011} & (TBD) & & \\ \cmidrule{1-4}
        \textbf{012} & (TBD) & & \\
        
        \bottomrule
    \end{tabular}
    
    \vspace{0.3cm}
    \small
%    \textit{Note}: \textbf{Phenomenon-Level Tests (P)} verify whether the psychological effect is present (critical for validation—must pass). \textbf{Data-Level Tests (D)} measure how closely the quantitative results match human baseline data using equivalence testing with tolerance $\delta = 0.2$ (non-critical—provides additional validation).
\end{table*}
